\chapter{Agent-Enhanced Incident Triage}
\label{ch:triage}

% Working draft - S. Vene, 2026-02-08

\section{Introduction}
\label{sec:triage-intro}

Incident response consumes disproportionate platform engineering resources. Root cause analysis (RCA) alone takes 40--60\% of incident time, while alert fatigue degrades response effectiveness. The DORA 2024 report found over 39\% of respondents have little or no trust in AI tools, yet reducing MTTR remains critical for organizational performance.

This chapter examines how AI agents can augment human responders in the triage and diagnosis phases of incident management---the phases where observability data (Chapter~\ref{ch:observability}) must be transformed into actionable understanding and response.

\section{The Triage Challenge}
\label{sec:triage-challenge}

\subsection{Current State Pain Points}
\label{subsec:pain-points}

\begin{table}[htbp]
\centering
\caption{Incident Triage Challenges}
\label{tab:triage-challenges}
\begin{tabular}{ll}
\toprule
\textbf{Challenge} & \textbf{Impact} \\
\midrule
Information overload & Responders must correlate logs, metrics, traces, change history \\
Expertise bottleneck & Senior engineers required for diagnosis, 24/7 coverage challenges \\
Documentation lag & Runbooks incomplete or outdated \\
Coordination overhead & Multi-team incidents require extensive communication \\
Learning loss & Incident knowledge doesn't transfer to prevention \\
\bottomrule
\end{tabular}
\end{table}

\subsection{What Agents Can Address}
\label{subsec:agent-capabilities}

Agents excel at:
\begin{itemize}
    \item \textbf{Speed:} Parallel analysis of multiple data sources
    \item \textbf{Memory:} Recall of similar past incidents and their resolutions
    \item \textbf{Correlation:} Connecting signals across systems and time
    \item \textbf{Consistency:} Same triage quality at 3am as at 3pm
\end{itemize}

Agents struggle with:
\begin{itemize}
    \item \textbf{Novel failures:} Patterns outside training distribution
    \item \textbf{Business context:} Understanding which services matter most right now
    \item \textbf{Human coordination:} Navigating organizational politics during major incidents
    \item \textbf{Judgment calls:} When to escalate, who to page, what to communicate
\end{itemize}

\section{Pattern: Automated Initial Triage}
\label{sec:auto-triage}

\subsection{Problem}

When an alert fires, someone must assess severity, identify the affected system, and route to the appropriate team. This initial triage is repetitive but requires context that's scattered across tools.

\subsection{Solution}

An agent that performs initial triage on incoming alerts: classification, severity assessment, context gathering, and routing recommendation.

\subsection{Implementation}

\textbf{Input processing:}
\begin{itemize}
    \item Alert metadata (source, severity, labels)
    \item Service ownership data
    \item Recent changes (deployments, config updates)
    \item Historical alerts for same service/metric
\end{itemize}

\textbf{Triage output:}

\begin{lstlisting}[language=yaml,caption={Automated Triage Output Example}]
alert: "api-gateway-high-error-rate"
received: "2026-02-08T14:32:00Z"

triage_assessment:
  severity: P2  # Elevated from P3 based on impact analysis
  severity_reasoning: |
    Error rate affects customer-facing checkout flow.
    Currently 12% of requests failing (threshold: 5%).
    Business hours, peak traffic period.

  affected_services:
    - api-gateway (primary)
    - checkout-service (downstream impact)

  owner_team: payments-platform
  on_call: "@jane.smith"

  context_gathered:
    recent_deployments:
      - "auth-service v2.3.2 deployed 2 min ago"
    recent_changes:
      - "none in affected services"
    similar_alerts_24h: 0

  routing_recommendation: |
    Page payments-platform on-call.
    Include auth-service team given recent deployment.

  confidence: 0.85
\end{lstlisting}

\subsection{Autonomy Level}

L3 (Monitored): Agent triages and routes. Human reviews escalation decisions.

\subsection{Risks}

\begin{itemize}
    \item Misclassification leading to wrong team engagement
    \item Severity underestimation missing critical issues
    \item Over-alerting if agent is too conservative
\end{itemize}

\section{Pattern: Root Cause Hypothesis Generation}
\label{sec:rca-hypothesis}

\subsection{Problem}

Root cause analysis requires correlating signals across systems, recalling similar past incidents, and forming hypotheses. This is cognitively demanding, especially at 3am.

\subsection{Solution}

An agent that analyzes incident signals and produces ranked hypotheses with supporting evidence.

\subsection{Implementation}

\textbf{Hypothesis generation process:}

\begin{enumerate}
    \item \textbf{Signal Collection}
    \begin{itemize}
        \item Query metrics for anomalies in time window
        \item Search logs for errors/warnings
        \item Check deployment/change timeline
        \item Query trace data for latency breakdowns
    \end{itemize}

    \item \textbf{Pattern Matching}
    \begin{itemize}
        \item Compare current signals to past incidents
        \item Identify common failure patterns
        \item Check for known issues in similar systems
    \end{itemize}

    \item \textbf{Hypothesis Formation}
    \begin{itemize}
        \item Generate candidate root causes
        \item Assign confidence based on evidence strength
        \item Rank by likelihood and impact
    \end{itemize}

    \item \textbf{Evidence Packaging}
    \begin{itemize}
        \item Link to specific logs, metrics, traces
        \item Show reasoning chain
        \item Highlight uncertainties
    \end{itemize}
\end{enumerate}

\textbf{Example output:}

\begin{lstlisting}[caption={Root Cause Hypothesis Example}]
INCIDENT: checkout-service timeout errors

HYPOTHESIS 1 (Confidence: 85%)
Root Cause: Auth-service deployment introduced N+1 query pattern

Evidence:
- auth-service deployment at 14:30 UTC (2 min before incident)
- auth-service /validate endpoint latency: 400ms -> 2.5s
- auth-db-primary connection pool: 50/50 exhausted
- Checkout-service depends on auth-service for validation
- Similar pattern in INC-2391 (resolved by rollback)

Suggested action: Rollback auth-service to v2.3.1

---

HYPOTHESIS 2 (Confidence: 35%)
Root Cause: Traffic spike exceeding capacity

Evidence:
- Request rate 20% above daily average
- But: capacity headroom should handle 50% spike
- Partially explains symptoms, doesn't explain database exhaustion

Suggested action: Scale checkout-service replicas (may help, not root fix)

---

HYPOTHESIS 3 (Confidence: 15%)
Root Cause: Upstream provider degradation

Evidence:
- payment-provider API latency elevated
- But: this affects different code path than observed errors

Suggested action: Monitor, likely not primary cause
\end{lstlisting}

\subsection{Autonomy Level}

L2 (Assisted): Agent proposes hypotheses. Human investigates and decides.

\subsection{Industry Implementation: HolmesGPT}

Robusta.dev's HolmesGPT implements this pattern in production. At KubeCon 2026, practitioners reported 80\% correct root cause diagnosis in under 2 minutes. Their architecture enforces read-only access with PR approval gates---the agent diagnoses but cannot directly remediate, maintaining the human-in-the-loop required for L2 autonomy.

Key prerequisites they identify: strong observability infrastructure, consistent incident logging, and a structured service catalog mapping applications to resources and documentation. This validates the thesis's position that observability patterns (Chapter~\ref{ch:observability}) are prerequisites for effective triage automation.

\section{Pattern: Similar Incident Retrieval}
\label{sec:similar-incident}

\subsection{Problem}

Institutional knowledge about past incidents exists in post-mortems, chat logs, and engineer memories. During an incident, this knowledge is hard to access quickly.

\subsection{Solution}

An agent that maintains a searchable index of past incidents and retrieves relevant cases during active incidents.

\subsection{Implementation}

\textbf{Incident knowledge base:}
\begin{itemize}
    \item Post-mortem documents
    \item Incident ticket history
    \item Chat transcripts from incident channels
    \item Runbook execution logs
    \item Resolution actions taken
\end{itemize}

\textbf{Retrieval triggers:}
\begin{itemize}
    \item New incident opened $\rightarrow$ search for similar
    \item Manual query (``have we seen this before?'')
    \item Hypothesis validation (``was this root cause seen before?'')
\end{itemize}

\textbf{Similarity matching:}
\begin{itemize}
    \item Service/component affected
    \item Error signatures (log patterns, error codes)
    \item Symptom profiles (latency, error rate, saturation)
    \item Temporal patterns (time of day, deployment correlation)
\end{itemize}

\textbf{Output:}

\begin{lstlisting}[caption={Similar Incident Retrieval Output}]
SIMILAR INCIDENTS for checkout-service timeout

1. INC-2391 (1 month ago) - 92% match
   Summary: Payment-service connection pool exhaustion
   Root cause: Missing connection timeout, leaked connections
   Resolution: Code fix + pool size increase
   Relevance: Same symptom pattern, similar database involvement

2. INC-2847 (2 weeks ago) - 71% match
   Summary: Auth-service high latency
   Root cause: Inefficient query in new feature
   Resolution: Query optimization
   Relevance: Same upstream service, different symptom

3. INC-2156 (3 months ago) - 65% match
   Summary: Checkout cascade failure
   Root cause: Kubernetes node failure
   Resolution: Pod rescheduling, node replacement
   Relevance: Same service, different root cause class
\end{lstlisting}

\subsection{Autonomy Level}

L2 (Assisted): Agent retrieves and summarizes. Human evaluates relevance.

\section{Human-Agent Collaboration in Diagnosis}
\label{sec:collaboration}

\subsection{The Collaboration Model}
\label{subsec:collaboration-model}

Effective incident triage combines agent strengths (speed, memory, correlation) with human strengths (judgment, context, novel reasoning):

\begin{lstlisting}[caption={Human-Agent Collaboration Flow}]
                Incident Occurs
                      |
                      v
+----------------------------------------------+
| AGENT: Initial Triage                        |
| - Classify, route, gather context            |
| - Generate hypotheses                        |
| - Retrieve similar incidents                 |
+----------------------------------------------+
                      |
                      v
+----------------------------------------------+
| HUMAN: Hypothesis Evaluation                 |
| - Assess agent hypotheses                    |
| - Add business context                       |
| - Direct additional investigation            |
+----------------------------------------------+
                      |
                      v
+----------------------------------------------+
| AGENT: Deep Dive Support                     |
| - Run specific queries on request            |
| - Gather additional evidence                 |
| - Draft communications                       |
+----------------------------------------------+
                      |
                      v
+----------------------------------------------+
| HUMAN: Decision and Action                   |
| - Confirm root cause                         |
| - Approve remediation                        |
| - Coordinate response                        |
+----------------------------------------------+
\end{lstlisting}

\subsection{Communication Patterns}
\label{subsec:communication-patterns}

\textbf{Agent to Human:}
\begin{itemize}
    \item Present findings as hypotheses, not facts
    \item Include confidence levels and evidence links
    \item Highlight uncertainties and gaps
    \item Suggest next steps without demanding action
\end{itemize}

\textbf{Human to Agent:}
\begin{itemize}
    \item Direct specific investigations (``check auth-service logs for X'')
    \item Provide context agent lacks (``we're in a deploy freeze'')
    \item Override agent assessments when appropriate
    \item Confirm or reject hypotheses
\end{itemize}

\subsection{Trust Building}
\label{subsec:trust-building-triage}

From Chapter~\ref{ch:trust-framework}'s framework: triage agents should start at L2 (Assisted) and earn higher autonomy through demonstrated accuracy:

\begin{table}[htbp]
\centering
\caption{Promotion Thresholds for Triage Agents}
\label{tab:triage-promotion}
\begin{tabular}{ll}
\toprule
\textbf{Metric} & \textbf{Threshold for L3} \\
\midrule
Hypothesis accuracy & $>$80\% correct in top 3 \\
Routing accuracy & $>$95\% correct team \\
False positive rate & $<$10\% \\
Time operating without major error & $>$30 days \\
\bottomrule
\end{tabular}
\end{table}

\section{Memory-Enhanced Incident Resolution}
\label{sec:memory-enhanced}

\subsection{Overview}
\label{subsec:memory-overview}

Incident resolution is perhaps the strongest use case for agent long-term memory. Senior SREs are valuable precisely because they remember: ``Last time the database did this, it was a connection pool leak after the Thursday deploy.''

This section explores how memory systems can enhance incident agents while maintaining trust, and examines the failure modes unique to high-stakes operational memory.

\subsection{Memory Architecture for Incidents}
\label{subsec:memory-architecture}

Following Hu et al. (2025) ``Memory in the Age of AI Agents,'' incident agents need three memory types:

\begin{table}[htbp]
\centering
\caption{Memory Types for Incident Agents}
\label{tab:memory-types}
\begin{tabular}{lll}
\toprule
\textbf{Memory Type} & \textbf{Incident Application} & \textbf{Example} \\
\midrule
Factual & System topology, runbooks, config & ``Database pool default is 100 connections'' \\
Experiential & Past incidents, what worked/failed & ``Similar error in \#3892 $\rightarrow$ fixed by rollback'' \\
Working & Current investigation state & ``Already checked network, testing DB next'' \\
\bottomrule
\end{tabular}
\end{table}

\subsection{Memory Patterns from Practice}
\label{subsec:memory-patterns}

\textbf{Similar Incident Retrieval:} Semantic search over postmortem corpus---``Have we seen this before?''

\textbf{Error-Triggered Recall:} Don't retrieve unless anomaly detected. Bash error $\rightarrow$ surface past fixes for that error.

\textbf{Token-Budgeted Injection:} Greedy knapsack with caps. Most relevant memories without context pollution.

\textbf{Feedback Loops:} Track which recommendations were adopted. Deprioritize unused suggestions.

\subsection{Trust Framework for Memory}
\label{subsec:memory-trust}

\begin{table}[htbp]
\centering
\caption{Trust Factors for Memory-Enhanced Agents}
\label{tab:memory-trust}
\begin{tabular}{lll}
\toprule
\textbf{Factor} & \textbf{Memory Application} & \textbf{Design Implication} \\
\midrule
Observability & What memories influenced this decision? & Provenance logging, retrieval explanation \\
Reversibility & Can we undo memory-based actions? & Soft recommendations before automation \\
Blast Radius & Impact of wrong memory? & Graduated autonomy by domain \\
Autonomy & How much does memory drive action? & Human-in-loop for high-stakes \\
\bottomrule
\end{tabular}
\end{table}

\subsubsection{Memory-Specific Trust: Confidence Decay}

Memory confidence should degrade:
\begin{equation}
\text{confidence} = \text{base\_confidence} \times \text{decay\_factor}(\text{age}, \text{domain\_volatility})
\end{equation}

Infrastructure memories decay fast (topology changes). Runbook patterns decay slower. Historical facts (incident \#3892 happened) don't decay.

\subsection{Failure Modes}
\label{subsec:memory-failures}

\subsubsection{Staleness}

Infrastructure changed but knowledge wasn't updated. The topology map is wrong. The escalation path goes to someone who left.

\textbf{Mitigation:} Temporal validity tags, confidence decay, automated validation against current state.

\subsubsection{Overfitting}

Pattern worked 5 times, agent becomes overconfident. But those 5 incidents had the same root cause. New incident with same symptom, different cause $\rightarrow$ wrong remediation applied confidently.

\textbf{Mitigation:} Track conditions of success, require diversity in validation, decay confidence in novel contexts.

\subsubsection{Pollution (Adversarial)}

AGENTPOISON (NeurIPS 2024) and PoisonedRAG (USENIX 2025) demonstrate attacks on agent memory:
\begin{itemize}
    \item Malicious postmortem suggests harmful runbook
    \item Poisoned similar-incident matching
\end{itemize}

\textbf{Mitigation:} Cryptographic signing of authoritative knowledge, provenance tracking, anomaly detection on memory updates.

\subsubsection{Dilution}

Too much retrieved context overwhelms signal. 10 ``similar'' incidents, none truly relevant. Token budget consumed by noise.

\textbf{Mitigation:} Token budgeting, relevance thresholds, hierarchical retrieval (summary first, detail on demand).

\subsubsection{Feedback Loop Pathologies}

Agent suggests wrong root cause $\rightarrow$ operator accepts (time pressure) $\rightarrow$ postmortem records as ``correct'' $\rightarrow$ future incidents get same wrong answer $\rightarrow$ self-reinforcing error.

\textbf{Mitigation:} Distinguish ``operator accepted'' from ``operator validated.'' Track long-term outcomes. Periodic expert review.

\subsection{Trust Zones for Memory-Driven Actions}
\label{subsec:trust-zones}

\begin{table}[htbp]
\centering
\caption{Trust Zones for Memory-Driven Actions}
\label{tab:trust-zones}
\begin{tabular}{llll}
\toprule
\textbf{Zone} & \textbf{Memory Use} & \textbf{Human Role} & \textbf{Example} \\
\midrule
Advisory & Surface similar incidents & Full review & ``Consider checking pool based on \#3892'' \\
Guided & Pre-fill diagnostic commands & Approve/execute & ``Run these queries? [Execute/Skip]'' \\
Supervised & Execute known-safe actions & Exception handling & Auto-route pages, assign severity \\
Autonomous & Full remediation from memory & Post-hoc audit & Auto-restart for recognized patterns \\
\bottomrule
\end{tabular}
\end{table}

\textbf{Progression rule:} Agents earn autonomy through demonstrated accuracy, tracked per domain.

\section{Risks and Mitigations}
\label{sec:triage-risks}

\begin{table}[htbp]
\centering
\caption{Triage Pattern Risks and Mitigations}
\label{tab:triage-risks}
\begin{tabular}{ll}
\toprule
\textbf{Risk} & \textbf{Mitigation} \\
\midrule
Wrong diagnosis with high confidence & Present as hypothesis, require evidence review \\
Over-reliance---team skills atrophy & Periodic agent-free exercises, training \\
Coordination confusion & Clear role definitions, explicit handoffs \\
Sensitive data in communications & Redaction rules, output filtering \\
Alert storm amplification & Rate limiting, correlation before routing \\
\bottomrule
\end{tabular}
\end{table}

\section{Summary}
\label{sec:triage-summary}

Agent-enhanced triage transforms incident response from a purely human cognitive task to a human-agent collaboration:

\begin{table}[htbp]
\centering
\caption{Triage Pattern Summary}
\label{tab:triage-summary}
\begin{tabular}{lll}
\toprule
\textbf{Pattern} & \textbf{Problem Solved} & \textbf{Autonomy Level} \\
\midrule
Automated Initial Triage & Repetitive classification and routing & L3 \\
Root Cause Hypothesis & Correlation and pattern matching & L2 \\
Similar Incident Retrieval & Institutional knowledge access & L2 \\
Memory-Enhanced Resolution & Long-term operational learning & L2--L4 \\
\bottomrule
\end{tabular}
\end{table}

The next chapter examines the more sensitive question: when and how can agents participate in remediation?

\textit{This chapter addresses RQ2: To what extent can agentic reasoning improve incident triage and root cause analysis compared to existing SRE practices?}
